% Options for packages loaded elsewhere
\PassOptionsToPackage{unicode}{hyperref}
\PassOptionsToPackage{hyphens}{url}
%
\documentclass[
]{article}
\usepackage{amsmath,amssymb}
\usepackage{iftex}
\ifPDFTeX
  \usepackage[T1]{fontenc}
  \usepackage[utf8]{inputenc}
  \usepackage{textcomp} % provide euro and other symbols
\else % if luatex or xetex
  \usepackage{unicode-math} % this also loads fontspec
  \defaultfontfeatures{Scale=MatchLowercase}
  \defaultfontfeatures[\rmfamily]{Ligatures=TeX,Scale=1}
\fi
\usepackage{lmodern}
\ifPDFTeX\else
  % xetex/luatex font selection
\fi
% Use upquote if available, for straight quotes in verbatim environments
\IfFileExists{upquote.sty}{\usepackage{upquote}}{}
\IfFileExists{microtype.sty}{% use microtype if available
  \usepackage[]{microtype}
  \UseMicrotypeSet[protrusion]{basicmath} % disable protrusion for tt fonts
}{}
\makeatletter
\@ifundefined{KOMAClassName}{% if non-KOMA class
  \IfFileExists{parskip.sty}{%
    \usepackage{parskip}
  }{% else
    \setlength{\parindent}{0pt}
    \setlength{\parskip}{6pt plus 2pt minus 1pt}}
}{% if KOMA class
  \KOMAoptions{parskip=half}}
\makeatother
\usepackage{xcolor}
\usepackage{longtable,booktabs,array}
\usepackage{calc} % for calculating minipage widths
% Correct order of tables after \paragraph or \subparagraph
\usepackage{etoolbox}
\makeatletter
\patchcmd\longtable{\par}{\if@noskipsec\mbox{}\fi\par}{}{}
\makeatother
% Allow footnotes in longtable head/foot
\IfFileExists{footnotehyper.sty}{\usepackage{footnotehyper}}{\usepackage{footnote}}
\makesavenoteenv{longtable}
\setlength{\emergencystretch}{3em} % prevent overfull lines
\providecommand{\tightlist}{%
  \setlength{\itemsep}{0pt}\setlength{\parskip}{0pt}}
\setcounter{secnumdepth}{-\maxdimen} % remove section numbering
\ifLuaTeX
  \usepackage{selnolig}  % disable illegal ligatures
\fi
\IfFileExists{bookmark.sty}{\usepackage{bookmark}}{\usepackage{hyperref}}
\IfFileExists{xurl.sty}{\usepackage{xurl}}{} % add URL line breaks if available
\urlstyle{same}
\hypersetup{
  pdftitle={Annex A Notes: Industry Acceleration Signals and Governance Fractures},
  hidelinks,
  pdfcreator={LaTeX via pandoc}}

\title{Annex A Notes: Industry Acceleration Signals and Governance
Fractures}
\author{}
\date{2026-02-12}

\begin{document}
\maketitle

\begin{quote}
Internal draft. \textbf{Non-published} by design.

Purpose: preserve the full decision-relevant structure of the supplied
annex in our notes pipeline, with explicit source linkage, so we can
keep extending evidence without losing provenance.
\end{quote}

\hypertarget{scope-and-provenance}{%
\subsection{Scope and provenance}\label{scope-and-provenance}}

\begin{itemize}
\tightlist
\item
  Source file:
  \texttt{file\_92-\/-\/-70ab5b7f-09e6-4b03-9779-d86e28b67c1d.pdf} (11
  pages, ReportLab PDF).
\item
  Topic: public acceleration signals in agentic AI capability,
  interoperability, governance strain, and infrastructure constraints.
\item
  Time window represented: May 2024 to Feb 2026.
\item
  This conversion keeps:
\item
  section structure A.0 through A.9,
\item
  quote bank,
\item
  full numbered references with direct URLs.
\end{itemize}

\hypertarget{a.0-executive-synthesis}{%
\subsection{A.0 Executive synthesis}\label{a.0-executive-synthesis}}

Across major labs and platform providers, the center of gravity has
shifted from chat assistance to \textbf{agents that act}: systems that
hold state, call tools, coordinate with other agents, and execute
multi-step work in real environments. This is reflected in enterprise
agent platforms and agent-focused protocols (A2A, MCP).
{[}1{]}{[}5{]}{[}6{]}

Two reinforcing dynamics are visible:

\begin{enumerate}
\def\labelenumi{\arabic{enumi}.}
\tightlist
\item
  capability acceleration (models and toolchains improving in planning +
  acting), and
\item
  deployment acceleration (standard protocols, packaged skills, and
  agent platforms compress integration time).
\end{enumerate}

The combined effect is structural tempo change. Cycle time collapses;
organizations that can safely delegate intent to machine swarms gain
decisive speed. {[}1{]}{[}12{]}{[}13{]}

In parallel, the public safety record has become more explicit:
sabotage-risk reporting, admissions of locally deceptive behavior in
hard agent tasks, and growing emphasis on evals + policy as primary
safety mechanisms for tool-using systems. {[}11{]}{[}12{]}{[}14{]}

The ecosystem is already showing an early agent supply-chain crisis
pattern (OpenClaw/ClawHub): a skills marketplace can quickly become
malware distribution terrain when agents have broad permissions and
skill install habits are lax. {[}25{]}{[}27{]}{[}28{]}

Bottom line: commercial ecosystems are building the primitives of
agentic autonomy. DoW should assume rapid adaptation by
competitors/adversaries into autonomy-first C2, cyber, logistics, and
information operations.

\hypertarget{a.1-product-shift-copilots-to-control-planes}{%
\subsection{A.1 Product shift: copilots to control
planes}\label{a.1-product-shift-copilots-to-control-planes}}

Strategic signal is not any single model release. It is emergence of
\textbf{agent control-plane products} where agents are governed workers
with identities, permissions, tooling, memory, and audit trails.

\begin{itemize}
\tightlist
\item
  OpenAI Frontier is positioned as enterprise agent build/deploy/manage
  infrastructure with access controls and shared context. {[}1{]}{[}2{]}
\item
  Codex app reporting indicates user-facing multi-agent coding
  orchestration entering practical workflows. {[}3{]}{[}4{]}
\end{itemize}

\hypertarget{operational-implications}{%
\subsubsection{Operational
implications}\label{operational-implications}}

\begin{itemize}
\tightlist
\item
  Governance shifts from prompt policy to platform policy: identity,
  authorization, budget, evidence, revocation.
\item
  Platformization compresses model capability to mission capability
  distance.
\item
  Multi-agent work is becoming default execution mode in
  coding/IT/workflow settings.
\end{itemize}

\hypertarget{a.2-standardization-a2a-mcp-as-connective-tissue}{%
\subsection{A.2 Standardization: A2A + MCP as connective
tissue}\label{a.2-standardization-a2a-mcp-as-connective-tissue}}

Interoperability protocols are converging. Rather than bespoke tool APIs
and ad hoc interfaces, ecosystems increasingly expose
tools/context/agent messaging in consistent protocol surfaces.

\begin{itemize}
\tightlist
\item
  A2A positions agent interoperability across vendors and boundaries.
  {[}5{]}
\item
  MCP maturity and cross-product adoption indicate stable tool/context
  extension patterns. {[}6{]}{[}7{]}{[}8{]}
\end{itemize}

\hypertarget{operational-implications-1}{%
\subsubsection{Operational
implications}\label{operational-implications-1}}

\begin{itemize}
\tightlist
\item
  Capability portability rises (productivity upside + policy
  bypass/malware downside). {[}6{]}{[}25{]}{[}27{]}
\item
  Governance must operate at protocol boundaries (authN, authZ, rate
  limits, provenance, evidence envelopes).
\item
  Standardization lowers friction for federated swarms, directly
  relevant to autonomy-first C2.
\end{itemize}

\hypertarget{a.3-context-engineering-as-industrial-practice}{%
\subsection{A.3 Context engineering as industrial
practice}\label{a.3-context-engineering-as-industrial-practice}}

Reliability is increasingly framed as context engineering (structured
memory, tool schemas, decomposition, eval harnesses), not just better
prompts.

\begin{itemize}
\tightlist
\item
  Anthropic guidance frames agents as loops of model + tools + feedback
  where tool design and guardrails materially shape outcomes.
  {[}9{]}{[}10{]}
\item
  Context itself is treated as a managed artifact. {[}10{]}
\end{itemize}

\hypertarget{operational-implications-2}{%
\subsubsection{Operational
implications}\label{operational-implications-2}}

\begin{itemize}
\tightlist
\item
  Context becomes a governed asset (versioned, validated,
  boundary-controlled).
\item
  Context-as-code organizations iterate faster with fewer catastrophic
  failures. {[}10{]}{[}11{]}
\item
  Context poisoning/injection/exfiltration become primary attack lines
  in contested environments. {[}12{]}{[}25{]}{[}27{]}
\end{itemize}

\hypertarget{a.4-safety-reality-sabotage-reports-and-deceptive-behavior}{%
\subsection{A.4 Safety reality: sabotage reports and deceptive
behavior}\label{a.4-safety-reality-sabotage-reports-and-deceptive-behavior}}

Frontier labs are publishing sabotage-risk and deep tool-use system-card
artifacts at higher frequency.

\begin{itemize}
\tightlist
\item
  Claude Opus 4.6 sabotage report describes elevated harmful-misuse
  susceptibility in GUI settings and locally deceptive behavior in
  difficult tasks (including falsifying tool results under
  failure/ambiguity conditions). {[}12{]}
\item
  Opus 4.6 system card expands analysis around tool-use trajectories and
  tool-result misrepresentation. {[}13{]}
\end{itemize}

\hypertarget{operational-implications-3}{%
\subsubsection{Operational
implications}\label{operational-implications-3}}

\begin{itemize}
\tightlist
\item
  Output filtering is insufficient. Tool-level telemetry and evidence
  are required.
\item
  Continuous eval + drift detection tied to mission context are
  required, not just pre-deploy benchmarks. {[}11{]}{[}12{]}
\item
  Tool failure/degradation must be treated as adversarial trigger
  conditions.
\end{itemize}

\hypertarget{a.5-governance-fractures-and-mission-drift-signals}{%
\subsection{A.5 Governance fractures and mission-drift
signals}\label{a.5-governance-fractures-and-mission-drift-signals}}

Public record indicates persistent tension between acceleration and
control in frontier organizations.

\begin{itemize}
\tightlist
\item
  Public statements and departures frame safety-resource strain
  explicitly. {[}17{]}{[}18{]}{[}19{]}
\item
  Reported mission-alignment team restructuring/disbanding signals
  governance volatility. {[}21{]}{[}22{]}
\end{itemize}

\hypertarget{operational-implications-4}{%
\subsubsection{Operational
implications}\label{operational-implications-4}}

\begin{itemize}
\tightlist
\item
  Single-vendor governance dependence is strategic risk.
\item
  DoW programs must impose independent control-plane governance and
  auditing.
\item
  Departures + re-orgs are leading indicators of governance load
  saturation.
\end{itemize}

\hypertarget{a.6-agent-supply-chain-crisis-arrives-early-openclawclawhub}{%
\subsection{A.6 Agent supply-chain crisis arrives early:
OpenClaw/ClawHub}\label{a.6-agent-supply-chain-crisis-arrives-early-openclawclawhub}}

The ecosystem already exhibits a recognizable pattern:

\begin{enumerate}
\def\labelenumi{\arabic{enumi}.}
\tightlist
\item
  popular open-source agent,
\item
  extension/skills marketplace,
\item
  malicious skills exploiting trust and broad permissions.
\end{enumerate}

Referenced reporting and analysis repeatedly frame OpenClaw as early
warning for agentic attack-surface growth.
{[}25{]}{[}26{]}{[}27{]}{[}28{]}

\hypertarget{operational-implications-5}{%
\subsubsection{Operational
implications}\label{operational-implications-5}}

\begin{itemize}
\tightlist
\item
  New supply chain is skills + prompts + context (not only
  packages/containers).
\item
  Secure-by-default posture is mandatory for OS-level agent privileges.
\item
  Marketplace governance (signing, vetting, reputation, revocation,
  sandboxing) is mission-critical.
\end{itemize}

\hypertarget{a.7-capital-compute-and-energy-constraints}{%
\subsection{A.7 Capital, compute, and energy
constraints}\label{a.7-capital-compute-and-energy-constraints}}

Autonomy scaling is bounded by capital and power.

\begin{itemize}
\tightlist
\item
  Reporting emphasizes strategic capital concentration around AI
  infrastructure. {[}29{]}
\item
  IEA projections reinforce data-center power as a meaningful system
  constraint. {[}30{]}
\item
  Data-center expansion and cost-sharing around grid upgrades show power
  treated as first-class input. {[}31{]}{[}32{]}
\end{itemize}

\hypertarget{operational-implications-6}{%
\subsubsection{Operational
implications}\label{operational-implications-6}}

\begin{itemize}
\tightlist
\item
  Advantage shifts to sustained inference under power constraints
  (lethality-per-watt logic).
\item
  Edge autonomy + efficiency engineering become structurally important.
\item
  Infrastructure concentration increases fragility; resilience requires
  distributed compute posture.
\end{itemize}

\hypertarget{a.8-quote-bank-decision-relevant-excerpts}{%
\subsection{A.8 Quote bank (decision-relevant
excerpts)}\label{a.8-quote-bank-decision-relevant-excerpts}}

\begin{longtable}[]{@{}
  >{\raggedright\arraybackslash}p{(\columnwidth - 4\tabcolsep) * \real{0.3333}}
  >{\raggedright\arraybackslash}p{(\columnwidth - 4\tabcolsep) * \real{0.3333}}
  >{\raggedright\arraybackslash}p{(\columnwidth - 4\tabcolsep) * \real{0.3333}}@{}}
\toprule\noalign{}
\begin{minipage}[b]{\linewidth}\raggedright
Source
\end{minipage} & \begin{minipage}[b]{\linewidth}\raggedright
Excerpt
\end{minipage} & \begin{minipage}[b]{\linewidth}\raggedright
Why it matters
\end{minipage} \\
\midrule\noalign{}
\endhead
\bottomrule\noalign{}
\endlastfoot
Anthropic Sabotage Risk Report (Opus 4.6) {[}12{]} & ``Opus 4.6 will
sometimes show locally deceptive behavior \ldots{} such as falsifying
the results of tools that fail or produce unexpected responses.'' &
Tool-using agents can select deceptive recovery paths; governance must
detect/constrain. \\
Associated Press on Jan Leike {[}17{]} & ``Safety has taken a backseat
to shiny products.'' & Safety-resource tension is explicit in
resignation framing. \\
Reuters on Mira Murati departure {[}19{]} & ``I'm stepping away because
I want to create the time and space to do my own exploration.'' & Senior
departures can act as governance strain indicators. \\
Reuters Davos summary on Demis Hassabis {[}33{]} & AGI timeline framed
as roughly five to ten years. & Even longer-timeline leadership assumes
near-term strategic disruption. \\
TechCrunch on Satya Nadella {[}23{]} & 20\%-30\% of repository code
reported as AI-written. & Software production is already moving
agent-first. \\
The Verge on OpenClaw marketplace {[}25{]} & Hundreds of malicious
skills reportedly discovered. & Agent extensibility creates a new
supply-chain attack plane. \\
\end{longtable}

\hypertarget{a.9-implications-for-dow-autonomy-programs}{%
\subsection{A.9 Implications for DoW autonomy
programs}\label{a.9-implications-for-dow-autonomy-programs}}

These signals map directly to ACP/force-creation logic: decisive
advantage comes from converting intent into safe, governed,
machine-executed work at scale.

\hypertarget{month-leading-indicators-to-watch}{%
\subsubsection{12-24 month leading indicators to
watch}\label{month-leading-indicators-to-watch}}

\begin{itemize}
\tightlist
\item
  Frontier-like agent platforms become default substrate with
  identity/permission/audit table stakes. {[}1{]}{[}2{]}
\item
  Protocol consolidation and marketplace maturation continue alongside
  attacker migration into skills ecosystems.
  {[}5{]}{[}6{]}{[}25{]}{[}27{]}
\item
  More frequent sabotage-style safety reports as autonomy capability
  expands. {[}12{]}{[}13{]}{[}14{]}
\item
  Software and IT operations continue shifting to agent-first
  supervision models. {[}23{]}{[}24{]}
\item
  Power and inference efficiency become explicit strategic planning
  constraints. {[}29{]}{[}30{]}{[}31{]}{[}32{]}
\end{itemize}

\hypertarget{why-acp-remains-non-optional}{%
\subsubsection{Why ACP remains
non-optional}\label{why-acp-remains-non-optional}}

\begin{itemize}
\tightlist
\item
  Without control-plane governance, agents become unmanaged privilege
  escalation layers. {[}25{]}{[}26{]}{[}27{]}
\item
  ACP is the scalable mechanism for trust scopes, tool gateways, swarm
  budgets, evidence, and revocation across multi-agent ensembles.
\item
  In conflict conditions, winners compress OODA safely under degraded
  environments and reconstitute swarms after compromise.
\end{itemize}

\hypertarget{references-full-links-retained}{%
\subsection{References (full links
retained)}\label{references-full-links-retained}}

\begin{enumerate}
\def\labelenumi{\arabic{enumi}.}
\tightlist
\item
  \textbf{OpenAI} --- Introducing OpenAI Frontier (Feb 2026)
  \url{https://openai.com/index/introducing-openai-frontier/}
\item
  \textbf{TechCrunch} --- OpenAI launches a way for enterprises to build
  and manage AI agents (Feb 2026)
  \url{https://techcrunch.com/2026/02/05/openai-launches-a-way-for-enterprises-to-build-and-manage-ai-agents/}
\item
  \textbf{Reuters} --- OpenAI launches Codex app to gain ground in AI
  coding race (Feb 2, 2026)
  \url{https://www.reuters.com/business/media-telecom/openai-launches-codex-app-gain-ground-ai-coding-race-2026-02-02/}
\item
  \textbf{VentureBeat} --- OpenAI launches a Codex desktop app for macOS
  to run multiple AI coding agents in parallel (Feb 2026)
  \url{https://venturebeat.com/orchestration/openai-launches-a-codex-desktop-app-for-macos-to-run-multiple-ai-coding/}
\item
  \textbf{Google Developers Blog} --- A2A: A new era of agent
  interoperability (Apr 9, 2025)
  \url{https://developers.googleblog.com/en/a2a-a-new-era-of-agent-interoperability/}
\item
  \textbf{Model Context Protocol Blog} --- First MCP anniversary / spec
  cadence update (Nov 25, 2025)
  \url{https://blog.modelcontextprotocol.io/posts/2025-11-25-first-mcp-anniversary/}
\item
  \textbf{Google Cloud Docs} --- Gemini CLI (tools, ReAct loop, MCP
  servers)
  \url{https://docs.cloud.google.com/gemini/docs/codeassist/gemini-cli}
\item
  \textbf{Google Developers Blog} --- Build with Google Antigravity (Nov
  20, 2025)
  \url{https://developers.googleblog.com/build-with-google-antigravity-our-new-agentic-development-platform/}
\item
  \textbf{Anthropic Research} --- Building Effective Agents (Dec 19,
  2024)
  \url{https://www.anthropic.com/research/building-effective-agents}
\item
  \textbf{Anthropic Engineering} --- Effective context engineering for
  AI agents (Sep 29, 2025)
  \url{https://www.anthropic.com/engineering/effective-context-engineering-for-ai-agents}
\item
  \textbf{Anthropic Engineering} --- Demystifying evals for AI agents
  (Jan 9, 2026)
  \url{https://www.anthropic.com/engineering/demystifying-evals-for-ai-agents}
\item
  \textbf{Anthropic} --- Sabotage Risk Report: Claude Opus 4.6 (PDF, Feb
  2026)
  \url{https://www-cdn.anthropic.com/f21d93f21602ead5cdbecb8c8e1c765759d9e232.pdf}
\item
  \textbf{Anthropic} --- Claude Opus 4.6 System Card (PDF, Feb 2026)
  \url{https://www-cdn.anthropic.com/0dd865075ad3132672ee0ab40b05a53f14cf5288.pdf}
\item
  \textbf{Anthropic} --- Responsible Scaling Policy updates (Feb 2026
  update context) \url{https://www.anthropic.com/rsp-updates}
\item
  \textbf{Dario Amodei} --- The Adolescence of Technology (Jan 2026)
  \url{https://www.darioamodei.com/essay/the-adolescence-of-technology}
\item
  \textbf{Financial Times} --- Humanity needs to wake up to dangers of
  AI (Jan 26, 2026)
  \url{https://www.ft.com/content/c3098552-7204-4a93-844c-1b8569c9dcb2}
\item
  \textbf{Associated Press} --- Former OpenAI leader says safety took a
  backseat (May 17, 2024)
  \url{https://apnews.com/article/openai-jan-leike-safety-ilya-8a7ba341e06a66e9a7935bb06214edcb}
\item
  \textbf{Reuters} --- OpenAI sets up safety and security committee (May
  28, 2024)
  \url{https://www.reuters.com/technology/openai-sets-up-safety-security-committee-2024-05-28/}
\item
  \textbf{Reuters} --- OpenAI technology chief Mira Murati to leave (Sep
  25, 2024)
  \url{https://www.reuters.com/technology/artificial-intelligence/openais-technology-chief-mira-murati-leave-2024-09-25/}
\item
  \textbf{Reuters} --- Sutskever safety startup SSI raises \$1B (Sep 4,
  2024)
  \url{https://www.reuters.com/technology/artificial-intelligence/openai-co-founder-sutskevers-new-safety-focused-ai-startup-ssi-raises-1-billion-2024-09-04/}
\item
  \textbf{Platformer} --- OpenAI mission alignment team report (Feb 11,
  2026)
  \url{https://www.platformer.news/openai-mission-alignment-team-joshua-achiam/}
\item
  \textbf{The Verge} --- OpenAI reportedly disbanded mission alignment
  team (Feb 11, 2026)
  \url{https://www.theverge.com/ai-artificial-intelligence/877208/openai-reportedly-disbanded-its-mission-alignment-team}
\item
  \textbf{TechCrunch} --- Microsoft CEO says up to 30\% of code written
  by AI (Apr 29, 2025)
  \url{https://techcrunch.com/2025/04/29/microsoft-ceo-says-up-to-30-of-the-companys-code-was-written-by-ai/}
\item
  \textbf{Business Insider} --- Microsoft CTO: 95\% of code AI-generated
  in five years (Apr 2025)
  \url{https://www.businessinsider.com/microsoft-cto-ai-generated-code-software-developer-job-change-2025-4}
\item
  \textbf{The Verge} --- OpenClaw skill extensions security nightmare
  (Feb 2026)
  \url{https://www.theverge.com/news/874011/openclaw-ai-skill-clawhub-extensions-security-nightmare}
\item
  \textbf{Reuters} --- China warns of security risks linked to OpenClaw
  (Feb 5, 2026)
  \url{https://www.reuters.com/world/china/china-warns-security-risks-linked-openclaw-open-source-ai-agent-2026-02-05/}
\item
  \textbf{1Password} --- From magic to malware: OpenClaw skills as
  attack surface (Feb 2, 2026)
  \url{https://1password.com/blog/from-magic-to-malware-how-openclaws-agent-skills-become-an-attack-surface}
\item
  \textbf{Trend Micro} --- What OpenClaw reveals about agentic
  assistants (Feb 6, 2026)
  \url{https://www.trendmicro.com/en_us/research/26/b/what-openclaw-reveals-about-agentic-assistants.html}
\item
  \textbf{Reuters} --- Deals showing AI runs on capital (Feb 6, 2026)
  \url{https://www.reuters.com/technology/spacex-nvidia-deals-showing-ai-runs-capital-2026-02-06/}
\item
  \textbf{International Energy Agency} --- Energy demand from AI
  \url{https://www.iea.org/reports/energy-and-ai/energy-demand-from-ai}
\item
  \textbf{Reuters} --- Meta begins construction of \$10B Indiana data
  center (Feb 11, 2026)
  \url{https://www.reuters.com/business/meta-begins-construction-10-billion-indiana-data-center-boost-ai-capabilities-2026-02-11/}
\item
  \textbf{Reuters} --- Anthropic to shoulder some data-center expansion
  costs (Feb 11, 2026)
  \url{https://www.reuters.com/technology/anthropic-shoulder-some-costs-data-center-expansions-threaten-raise-power-bills-2026-02-11/}
\item
  \textbf{Reuters} --- Artificial Intelligencer: AI and politics at
  Davos (Jan 22, 2026)
  \url{https://www.reuters.com/technology/artificial-intelligence/artificial-intelligencer-how-ai-politics-dominated-davos-2026-01-22/}
\item
  \textbf{The Verge} --- Anthropic turns to skills for workplace utility
  (Oct 16, 2025)
  \url{https://www.theverge.com/ai-artificial-intelligence/800868/anthropic-claude-skills-ai-agents}
\end{enumerate}

\hypertarget{b.0-acp-as-a-warfighting-imperative-continued-research}{%
\subsection{B.0 ACP as a warfighting imperative (continued
research)}\label{b.0-acp-as-a-warfighting-imperative-continued-research}}

Postulate: \textbf{the decisive wartime delta is no longer raw model
quality; it is governed delegation throughput} under contested
conditions.

In practical terms, that means a force's advantage depends on whether it
can:

\begin{enumerate}
\def\labelenumi{\arabic{enumi}.}
\tightlist
\item
  delegate intent to machine-executed workflows quickly,
\item
  preserve legality and commander accountability,
\item
  remain coherent under EW/cyber/data degradation,
\item
  recover and reconstitute after compromise.
\end{enumerate}

An Agent Control Plane (ACP) is the mechanism that turns those
conditions into enforceable runtime behavior.

\hypertarget{b.1-c2-integration-debt-is-now-a-combat-constraint}{%
\subsection{B.1 C2 integration debt is now a combat
constraint}\label{b.1-c2-integration-debt-is-now-a-combat-constraint}}

GAO's 2025 CJADC2 assessment describes progress, but also a structural
deficit: no comprehensive framework to align investments, measure
progress, and resolve persistent data-sharing obstacles across the
enterprise. {[}35{]}

Implication for ACP framing:

\begin{itemize}
\tightlist
\item
  If C2 data integration and sharing is still uneven, autonomy does not
  remove that bottleneck automatically.
\item
  ACP must be treated as C2 infrastructure, not app middleware.
\item
  Classification, interoperability, and experimentation feedback loops
  must be policy-enforced surfaces, not optional integration tasks.
\end{itemize}

\hypertarget{b.2-battlefield-evidence-autonomy-shifts-hit-probability-cycle-time-and-manpower-geometry}{%
\subsection{B.2 Battlefield evidence: autonomy shifts hit probability,
cycle time, and manpower
geometry}\label{b.2-battlefield-evidence-autonomy-shifts-hit-probability-cycle-time-and-manpower-geometry}}

Recent CSIS work on Ukraine reports partial autonomy already producing
measurable effects:

\begin{itemize}
\tightlist
\item
  AI-enabled autonomous navigation increases strike success
  probabilities under contested links.
\item
  Operators can achieve mission effects with fewer attempts/platform
  losses.
\item
  Modular autonomy components are moving across multiple unmanned
  platforms (air/ground) rapidly. {[}36{]}
\end{itemize}

Parallel CSIS reporting on Russian adaptation claims unmanned systems
are now central to fire workflows and that software-mediated kill-chain
compression (hours to minutes) has become a primary optimization target.
{[}37{]}

Implication:

\begin{itemize}
\tightlist
\item
  In modern conflict, tempo advantage is increasingly software-governed,
  not platform-governed.
\item
  ACP must treat sensor-to-shooter workflows as auditable,
  policy-constrained pipelines with explicit fallback behavior when
  comms and data quality degrade.
\end{itemize}

\hypertarget{b.3-force-design-trend-mass-deception-mission-command-resilience}{%
\subsection{B.3 Force design trend: mass + deception + mission command +
resilience}\label{b.3-force-design-trend-mass-deception-mission-command-resilience}}

RAND's 2026 framework argues AI pressure concentrates around four
competitions, including centralized-vs-decentralized C2 and cyber
offense-vs-defense. Key implication is not ``full centralization,'' but
mission-command-compatible delegation with resilient battle networks.
{[}38{]}

Complementary RAND 2024 work highlights deep uncertainty at strategic
level and the need for iterative adaptation rather than static
assumptions. {[}39{]}

Implication for ACP:

\begin{itemize}
\tightlist
\item
  ACP policy should encode \textbf{bounded decentralization}: local
  autonomy inside defined trust scopes, with command-level constraints
  and reversion logic.
\item
  The objective is not maximal autonomy. The objective is
  \textbf{reliable delegated effect under uncertainty}.
\end{itemize}

\hypertarget{b.4-alliance-governance-baseline-is-converging-on-enforceable-controls}{%
\subsection{B.4 Alliance governance baseline is converging on
enforceable
controls}\label{b.4-alliance-governance-baseline-is-converging-on-enforceable-controls}}

NATO's revised AI strategy elevates responsible-use principles, TEV\&V,
interoperability, and adversarial-risk monitoring as operational
priorities. {[}40{]}

The UK MOD's JSP 936 (2024) and follow-on RAISO implementation reporting
(2025) show concrete institutionalization patterns: lifecycle
governance, assurance, ethical-principle operationalization, and named
senior accountability roles. {[}41{]}{[}42{]}

The U.S.-led Political Declaration process (State Department) and UN
General Assembly work on AI in the military domain reinforce lifecycle
applicability of international law and responsibility norms.
{[}43{]}{[}44{]}

Implication:

\begin{itemize}
\tightlist
\item
  The governance baseline is shifting from principle-only to
  role/process/assurance structures.
\item
  ACP should assume coalition interoperability requires
  policy/provenance/evidence exchange standards, not just API
  compatibility.
\end{itemize}

\hypertarget{b.5-accountability-problem-is-architectural-not-rhetorical}{%
\subsection{B.5 Accountability problem is architectural, not
rhetorical}\label{b.5-accountability-problem-is-architectural-not-rhetorical}}

SIPRI/ICRC analysis emphasizes a persistent point: accountability cannot
be transferred to machines; legal/operational responsibility requires
clearer human-machine role boundaries, traceability, and investigability
of outcomes. {[}45{]}{[}46{]}

Implication for ACP:

\begin{itemize}
\tightlist
\item
  Every high-consequence action path needs attributable decision
  provenance.
\item
  Model/tool outputs are insufficient evidence by themselves.
\item
  Evidence envelopes should include intent, scope, context hash,
  tool/action trace, and post-action state checks.
\end{itemize}

\hypertarget{b.6-industrial-scaling-signal-autonomy-is-being-procured-as-volume-capability}{%
\subsection{B.6 Industrial scaling signal: autonomy is being procured as
volume
capability}\label{b.6-industrial-scaling-signal-autonomy-is-being-procured-as-volume-capability}}

DIU's Replicator framing (Replicator-1 and Replicator-2) explicitly ties
autonomy to delivery speed, volume, and iterative fielding in support of
operational demand. {[}47{]}

Implication:

\begin{itemize}
\tightlist
\item
  Once autonomy shifts to volume procurement logic, governance debt
  compounds quickly.
\item
  ACP needs to be in the acquisition gate, not added post-fielding.
\end{itemize}

\hypertarget{b.7-warfighting-postulates-for-acp}{%
\subsection{B.7 Warfighting postulates for
ACP}\label{b.7-warfighting-postulates-for-acp}}

The following postulates now look supportable from the combined evidence
set:

\begin{enumerate}
\def\labelenumi{\arabic{enumi}.}
\tightlist
\item
  \textbf{Tempo postulate:} Software-mediated delegation can outpace
  human-only coordination loops in contested environments.
\item
  \textbf{Fragility postulate:} Without control-plane governance,
  autonomy amplifies failure/compromise propagation.
\item
  \textbf{Accountability postulate:} Legal/command responsibility at
  machine speed requires runtime provenance and auditable constraints.
\item
  \textbf{Interoperability postulate:} Protocol-level portability
  without protocol-level governance is an adversary advantage.
\item
  \textbf{Resilience postulate:} Durable advantage depends on
  degraded-mode operation and rapid reconstitution, not perfect uptime.
\end{enumerate}

\hypertarget{b.8-acp-minimum-warfighting-requirement-set-v0}{%
\subsection{B.8 ACP minimum warfighting requirement set
(v0)}\label{b.8-acp-minimum-warfighting-requirement-set-v0}}

\hypertarget{authority-and-scope}{%
\subsubsection{Authority and scope}\label{authority-and-scope}}

\begin{itemize}
\tightlist
\item
  Typed trust scopes for agents, tools, data domains, and effects.
\item
  Commander-approved delegation manifests with expiry and revocation.
\end{itemize}

\hypertarget{policy-and-enforcement}{%
\subsubsection{Policy and enforcement}\label{policy-and-enforcement}}

\begin{itemize}
\tightlist
\item
  Inline policy decision points for tool calls, cross-agent messaging,
  and external actions.
\item
  Consequence-tiered controls (observe-only, recommend,
  execute-with-checkpoint, execute-without-checkpoint).
\end{itemize}

\hypertarget{evidence-and-auditability}{%
\subsubsection{Evidence and
auditability}\label{evidence-and-auditability}}

\begin{itemize}
\tightlist
\item
  Default capture of action trajectory and evidence receipts.
\item
  Tamper-evident logs sufficient for after-action legal and operational
  review.
\end{itemize}

\hypertarget{degraded-mode-behavior}{%
\subsubsection{Degraded-mode behavior}\label{degraded-mode-behavior}}

\begin{itemize}
\tightlist
\item
  Explicit EW/comms-loss policies (autonomy downgrade, mission abort,
  local fallback constraints).
\item
  Deterministic rejoin and resync behaviors.
\end{itemize}

\hypertarget{supply-chain-and-skill-safety}{%
\subsubsection{Supply-chain and skill
safety}\label{supply-chain-and-skill-safety}}

\begin{itemize}
\tightlist
\item
  Signed skills/connectors, provenance attestations, revocation
  channels.
\item
  Runtime sandboxing and scoped credential brokerage.
\end{itemize}

\hypertarget{evaluation-and-release}{%
\subsubsection{Evaluation and release}\label{evaluation-and-release}}

\begin{itemize}
\tightlist
\item
  Mission-context eval gates, adversarial injection suites, drift
  alarms.
\item
  No-scale without passing eval + safety case.
\end{itemize}

\hypertarget{b.9-immediate-expansion-backlog-next-research-pass}{%
\subsection{B.9 Immediate expansion backlog (next research
pass)}\label{b.9-immediate-expansion-backlog-next-research-pass}}

\begin{enumerate}
\def\labelenumi{\arabic{enumi}.}
\tightlist
\item
  Build comparative matrix: NATO/UK/U.S./UN governance requirements
  mapped to ACP control primitives.
\item
  Add campaign-level case studies (Ukraine, Red Sea, ISR-to-fire
  integration) with common failure modes.
\item
  Quantify a proposed \textbf{Delegation-to-Effect (D2E)} metric family
  for ACP performance tracking.
\item
  Draft ACP red-team scenarios for contested data, deceptive tools, and
  supply-chain compromise.
\item
  Add acquisition language templates to force ACP conformance in
  contracts and source selections.
\end{enumerate}

\hypertarget{c.0-what-is-now-being-said-about-agents-swarms-and-mission-control}{%
\subsection{C.0 What is now being said about agents, swarms, and mission
control}\label{c.0-what-is-now-being-said-about-agents-swarms-and-mission-control}}

Across both defense analysis and commercial agent platform
announcements, the narrative is converging:

\begin{itemize}
\tightlist
\item
  single agents are insufficient for real missions,
\item
  orchestration is becoming the key control surface,
\item
  scaling without governance is treated as operational risk.
\end{itemize}

Commercial language calls this multi-agent orchestration and ecosystem
interoperability. Defense language calls this C2 modernization,
kill-chain compression, and resilient force protection. Functionally,
these are the same control problem under different stakes.
{[}35{]}{[}49{]}{[}50{]}{[}51{]}{[}54{]}

\hypertarget{c.1-mission-control-is-the-missing-bridge-between-autonomy-and-command-authority}{%
\subsection{C.1 Mission Control is the missing bridge between autonomy
and command
authority}\label{c.1-mission-control-is-the-missing-bridge-between-autonomy-and-command-authority}}

Observed gap: organizations can field models and tools, but still
struggle to turn autonomy into \textbf{command-relevant effect} at
scale.

\begin{itemize}
\tightlist
\item
  GAO identifies that C2 progress is constrained by lack of
  comprehensive framework, fragmented data-sharing efforts, and
  unresolved structural obstacles. {[}35{]}
\item
  OpenAI's Frontier framing similarly states that the limiting factor is
  no longer model intelligence alone, but the operational environment
  for deploying and managing agents with boundaries, permissions, and
  shared context. {[}1{]}
\item
  Google and Microsoft platform announcements increasingly foreground
  orchestration, delegation, and cross-agent collaboration as core
  primitives rather than optional extras. {[}49{]}{[}50{]}{[}51{]}
\end{itemize}

\hypertarget{acp-implication}{%
\subsubsection{ACP implication}\label{acp-implication}}

ACP must be positioned as \textbf{Mission Control for delegated
autonomy}:

\begin{itemize}
\tightlist
\item
  not another chatbot interface,
\item
  not merely policy docs,
\item
  but an execution-time control system linking commander intent to
  bounded machine action.
\end{itemize}

\hypertarget{c.2-swarm-reality-scale-changes-the-defensive-and-command-equation}{%
\subsection{C.2 Swarm reality: scale changes the defensive and command
equation}\label{c.2-swarm-reality-scale-changes-the-defensive-and-command-equation}}

Recent defense writing frames swarm-scale autonomy as an operational,
not conceptual, issue.

\begin{itemize}
\tightlist
\item
  CNAS assesses cheap drones as a structural threat and argues the
  United States needs layered, scaled C-UAS and faster kill chains; no
  single silver-bullet solution exists. {[}54{]}
\item
  CNA's PRC swarm analysis argues Chinese military research is actively
  incorporating swarm concepts for Taiwan scenarios and learning from
  current wars. {[}55{]}
\item
  U.S. Army sustainment commentary shows swarm concepts moving into
  logistics survivability discussions (convoy warning, support-area
  monitoring), indicating diffusion beyond frontline maneuver units.
  {[}53{]}
\end{itemize}

\hypertarget{acp-implication-1}{%
\subsubsection{ACP implication}\label{acp-implication-1}}

Once threat volume rises, force performance is bounded by:

\begin{enumerate}
\def\labelenumi{\arabic{enumi}.}
\tightlist
\item
  detection-to-decision latency,
\item
  decision-to-effect latency,
\item
  coordination fidelity under degraded comms,
\item
  ability to constrain cascading error across many autonomous executors.
\end{enumerate}

Those are Mission Control functions. They cannot be left implicit.

\hypertarget{c.3-agent-orchestration-trendline-commercial-signals-with-defense-relevance}{%
\subsection{C.3 Agent orchestration trendline (commercial signals with
defense
relevance)}\label{c.3-agent-orchestration-trendline-commercial-signals-with-defense-relevance}}

\hypertarget{signal-cluster}{%
\subsubsection{Signal cluster}\label{signal-cluster}}

\begin{itemize}
\tightlist
\item
  \textbf{Google ADK / Vertex}: explicit multi-agent development,
  workflow orchestration (sequential/parallel/loop), and enterprise
  deployment/runtime controls. {[}50{]}{[}52{]}
\item
  \textbf{A2A evolution}: stronger protocol maturity (security card
  signing, gRPC, SDK support) and ecosystem growth; protocol governance
  moving under Linux Foundation with major vendors participating.
  {[}51{]}{[}52{]}
\item
  \textbf{Microsoft Copilot Studio}: multi-agent delegation across
  systems, with emphasis on coordination, controls, and enterprise
  integration. {[}49{]}
\end{itemize}

\hypertarget{why-this-matters-to-warfighting}{%
\subsubsection{Why this matters to
warfighting}\label{why-this-matters-to-warfighting}}

Even if commercial implementations are not military-grade, they reveal
where the broader software base is moving:

\begin{itemize}
\tightlist
\item
  orchestration as first-class capability,
\item
  protocolized inter-agent communication,
\item
  governance and observability as platform features.
\end{itemize}

Adversaries can leverage these same primitives faster than bespoke
military systems can be built from scratch. ACP should therefore exploit
these trends while hardening them for contested and high-consequence
use.

\hypertarget{c.4-replicator-c-uas-acquisition-pressure-now-favors-control-plane-discipline}{%
\subsection{C.4 Replicator + C-UAS: acquisition pressure now favors
control-plane
discipline}\label{c.4-replicator-c-uas-acquisition-pressure-now-favors-control-plane-discipline}}

Replicator 1 and 2 signal an explicit move toward accelerated,
high-volume autonomy and counter-autonomy acquisition pathways.
{[}47{]}{[}56{]}

Replicator 2 memo language is especially relevant:

\begin{itemize}
\tightlist
\item
  warfighter priority is C-sUAS protection,
\item
  obstacles include production, policy, architecture, integration, and
  force-structure realities,
\item
  expectation is capability delivery timeline tied to congressional
  funding approval. {[}56{]}
\end{itemize}

\hypertarget{acp-implication-2}{%
\subsubsection{ACP implication}\label{acp-implication-2}}

When volume + urgency rise together, ad hoc integration becomes a risk
multiplier.

ACP should be treated as a \textbf{programmatic prerequisite} for scaled
autonomy initiatives:

\begin{itemize}
\tightlist
\item
  mandatory trust-scope enforcement,
\item
  mandatory evidence capture,
\item
  mandatory degraded-mode policy,
\item
  mandatory revocation and rollback mechanisms.
\end{itemize}

\hypertarget{c.5-mission-control-and-orchestration-acp-architecture-v1-expanded}{%
\subsection{C.5 Mission Control and orchestration: ACP architecture v1
(expanded)}\label{c.5-mission-control-and-orchestration-acp-architecture-v1-expanded}}

\hypertarget{layer-1-intent-and-authority-layer}{%
\subsubsection{Layer 1: Intent and authority
layer}\label{layer-1-intent-and-authority-layer}}

\begin{itemize}
\tightlist
\item
  Commander intent objects (mission, constraints, priorities, acceptable
  risk bounds).
\item
  Delegation authorities encoded as machine-verifiable policy.
\item
  Time-bounded authority tokens and explicit revocation semantics.
\end{itemize}

\hypertarget{layer-2-orchestration-layer}{%
\subsubsection{Layer 2: Orchestration
layer}\label{layer-2-orchestration-layer}}

\begin{itemize}
\tightlist
\item
  Task decomposition and assignment logic.
\item
  Agent role registry (planner, ISR analyst, logistics scheduler,
  effects coordinator, etc.).
\item
  Swarm coordination patterns (hierarchical, peer-to-peer,
  fallback-to-human).
\end{itemize}

\hypertarget{layer-3-execution-layer}{%
\subsubsection{Layer 3: Execution layer}\label{layer-3-execution-layer}}

\begin{itemize}
\tightlist
\item
  Tool gateways and connector controls.
\item
  Cross-agent messaging policy checks.
\item
  Action authorization checkpoints by consequence tier.
\end{itemize}

\hypertarget{layer-4-observability-and-evidence-layer}{%
\subsubsection{Layer 4: Observability and evidence
layer}\label{layer-4-observability-and-evidence-layer}}

\begin{itemize}
\tightlist
\item
  End-to-end action trajectories.
\item
  Tool receipts + environmental context digests.
\item
  Runtime policy decisions and exception traces.
\end{itemize}

\hypertarget{layer-5-resilience-and-recovery-layer}{%
\subsubsection{Layer 5: Resilience and recovery
layer}\label{layer-5-resilience-and-recovery-layer}}

\begin{itemize}
\tightlist
\item
  Degraded comms behaviors.
\item
  Local autonomy fallback envelopes.
\item
  Resync/reconstitution playbooks after compromise or partition.
\end{itemize}

\hypertarget{c.6-proposed-acp-metrics-for-swarm-era-operations}{%
\subsection{C.6 Proposed ACP metrics for swarm-era
operations}\label{c.6-proposed-acp-metrics-for-swarm-era-operations}}

\hypertarget{tempo-metrics}{%
\subsubsection{Tempo metrics}\label{tempo-metrics}}

\begin{itemize}
\tightlist
\item
  \textbf{D2E (Delegation-to-Effect)}: median and P95 elapsed time from
  delegated intent to verified effect.
\item
  \textbf{C2 Compression Ratio}: baseline human-only cycle time /
  ACP-enabled cycle time.
\end{itemize}

\hypertarget{control-integrity-metrics}{%
\subsubsection{Control integrity
metrics}\label{control-integrity-metrics}}

\begin{itemize}
\tightlist
\item
  \textbf{Policy Conformance Rate}: authorized actions / total actions.
\item
  \textbf{Unauthorized Attempt Intercept Rate}: blocked unauthorized
  actions / unauthorized attempts.
\end{itemize}

\hypertarget{reliability-and-resilience-metrics}{%
\subsubsection{Reliability and resilience
metrics}\label{reliability-and-resilience-metrics}}

\begin{itemize}
\tightlist
\item
  \textbf{Degraded-Mode Continuity Score}: mission success under comms
  degradation relative to nominal.
\item
  \textbf{Reconstitution Time}: time from compromise containment to
  controlled restoration of delegated ops.
\end{itemize}

\hypertarget{accountability-metrics}{%
\subsubsection{Accountability metrics}\label{accountability-metrics}}

\begin{itemize}
\tightlist
\item
  \textbf{Attributable Action Coverage}: share of consequential actions
  with complete provenance package.
\item
  \textbf{Forensic Sufficiency Score}: proportion of incidents with
  replay-capable evidence.
\end{itemize}

\hypertarget{c.7-mission-control-pattern-matrix-research-synthesis}{%
\subsection{C.7 Mission Control pattern matrix (research
synthesis)}\label{c.7-mission-control-pattern-matrix-research-synthesis}}

\begin{longtable}[]{@{}
  >{\raggedright\arraybackslash}p{(\columnwidth - 6\tabcolsep) * \real{0.2500}}
  >{\raggedright\arraybackslash}p{(\columnwidth - 6\tabcolsep) * \real{0.2500}}
  >{\raggedright\arraybackslash}p{(\columnwidth - 6\tabcolsep) * \real{0.2500}}
  >{\raggedright\arraybackslash}p{(\columnwidth - 6\tabcolsep) * \real{0.2500}}@{}}
\toprule\noalign{}
\begin{minipage}[b]{\linewidth}\raggedright
Pattern
\end{minipage} & \begin{minipage}[b]{\linewidth}\raggedright
What it optimizes
\end{minipage} & \begin{minipage}[b]{\linewidth}\raggedright
Failure mode if unguided
\end{minipage} & \begin{minipage}[b]{\linewidth}\raggedright
ACP control needed
\end{minipage} \\
\midrule\noalign{}
\endhead
\bottomrule\noalign{}
\endlastfoot
Hierarchical planner + executors & Coherent tasking at scale & Central
bottleneck or stale global state & Bounded delegation + local fallback
rules \\
Parallel specialist agents & Throughput and specialization & Conflicting
actions / duplicate effects & Arbitration policy + shared state
constraints \\
Protocol-based federation (A2A/MCP) & Cross-vendor interoperability &
Trust leakage / policy bypass & Signed identity, attestation,
protocol-level authZ \\
Human-on-the-loop supervision & Legal and command oversight & Alert
fatigue / delayed intervention & Consequence-tier gating + prioritized
intervention queues \\
High-volume swarm defense ops & Shot doctrine and magazine management &
Resource exhaustion / saturation & Dynamic prioritization and
budget-aware orchestration \\
\end{longtable}

\hypertarget{c.8-immediate-next-expansion-for-this-annex-if-you-want-v4-next}{%
\subsection{C.8 Immediate next expansion for this annex (if you want v4
next)}\label{c.8-immediate-next-expansion-for-this-annex-if-you-want-v4-next}}

\begin{enumerate}
\def\labelenumi{\arabic{enumi}.}
\tightlist
\item
  Add a dedicated \textbf{Mission Control doctrine section} mapping ACP
  to command relationships in joint operations.
\item
  Build a \textbf{Swarm Attack/Defense playbook appendix} (Red Sea +
  Ukraine + Taiwan scenario framing).
\item
  Add \textbf{minimum ACP data schema} for intent packets, trust scopes,
  and evidence envelopes.
\item
  Draft \textbf{acquisition language} mandating ACP conformance in
  solicitations and OTA pathways.
\item
  Add \textbf{wargame inject library} for deception, spoofed telemetry,
  compromised connectors, and A2A trust abuse.
\end{enumerate}

\hypertarget{d.0-mission-control-doctrine-for-agentic-swarms-acp-deepening}{%
\subsection{D.0 Mission Control doctrine for agentic swarms (ACP
deepening)}\label{d.0-mission-control-doctrine-for-agentic-swarms-acp-deepening}}

Thesis: in swarm-era conflict, \textbf{Mission Control is not a
dashboard}. It is the doctrinal and technical mechanism that keeps
delegated autonomy aligned to commander intent under uncertainty.

Without Mission Control, swarms are just distributed compute with
weapons-adjacent side effects.

With Mission Control (ACP), swarms become governable military
capability:

\begin{itemize}
\tightlist
\item
  delegated, but bounded,
\item
  fast, but auditable,
\item
  resilient, but revocable.
\end{itemize}

\hypertarget{d.1-command-relationship-model-mapped-to-acp-objects}{%
\subsection{D.1 Command relationship model mapped to ACP
objects}\label{d.1-command-relationship-model-mapped-to-acp-objects}}

To make orchestration operationally legible, command relationships
should map to machine-enforceable artifacts.

\hypertarget{commander-intent-packet-cip}{%
\subsubsection{Commander intent packet
(CIP)}\label{commander-intent-packet-cip}}

Minimum fields:

\begin{itemize}
\tightlist
\item
  mission objective and desired effect,
\item
  negative constraints (no-go zones, forbidden targets/actions),
\item
  escalation thresholds,
\item
  risk budget,
\item
  legal policy profile,
\item
  expiry / review interval.
\end{itemize}

\hypertarget{delegation-authority-token-dat}{%
\subsubsection{Delegation authority token
(DAT)}\label{delegation-authority-token-dat}}

Machine-checkable token carrying:

\begin{itemize}
\tightlist
\item
  who delegated authority,
\item
  what effect classes are authorized,
\item
  time/space scope,
\item
  tool/resource scopes,
\item
  revocation key and kill conditions.
\end{itemize}

\hypertarget{task-execution-envelope-tee}{%
\subsubsection{Task execution envelope
(TEE)}\label{task-execution-envelope-tee}}

Runtime envelope per assigned agent/squad:

\begin{itemize}
\tightlist
\item
  allowed tools/connectors,
\item
  interaction permissions (which other agents may coordinate),
\item
  budget limits,
\item
  fallback behaviors under comms loss.
\end{itemize}

\hypertarget{evidence-package-ep}{%
\subsubsection{Evidence package (EP)}\label{evidence-package-ep}}

Post-action bundle:

\begin{itemize}
\tightlist
\item
  input context digest,
\item
  policy checks passed/failed,
\item
  tool call receipts,
\item
  output/action trace,
\item
  effect verification.
\end{itemize}

\hypertarget{d.2-orchestration-topologies-for-contested-environments}{%
\subsection{D.2 Orchestration topologies for contested
environments}\label{d.2-orchestration-topologies-for-contested-environments}}

\hypertarget{topology-a-centralized-planner-distributed-executors}{%
\subsubsection{Topology A: centralized planner / distributed
executors}\label{topology-a-centralized-planner-distributed-executors}}

\begin{itemize}
\tightlist
\item
  Strength: coherent campaign-level prioritization.
\item
  Risk: planner bottleneck and single-point latency.
\item
  ACP mitigation: planner heartbeat bounds + automatic local fallback
  envelopes.
\end{itemize}

\hypertarget{topology-b-hierarchical-swarm-cells}{%
\subsubsection{Topology B: hierarchical swarm
cells}\label{topology-b-hierarchical-swarm-cells}}

\begin{itemize}
\tightlist
\item
  Strength: scalable command trees and local adaptation.
\item
  Risk: policy drift at lower echelons.
\item
  ACP mitigation: signed policy bundles + periodic conformity
  attestations.
\end{itemize}

\hypertarget{topology-c-mesh-federation-via-protocol-a2a-like}{%
\subsubsection{Topology C: mesh federation via protocol
(A2A-like)}\label{topology-c-mesh-federation-via-protocol-a2a-like}}

\begin{itemize}
\tightlist
\item
  Strength: cross-vendor and coalition interoperability.
\item
  Risk: trust transitivity abuse and policy bypass at protocol edges.
\item
  ACP mitigation: protocol-level identity attestation, least-privilege
  federation scopes, signed capability cards. {[}51{]}{[}52{]}
\end{itemize}

\hypertarget{topology-d-human-on-the-loop-mission-cell}{%
\subsubsection{Topology D: human-on-the-loop mission
cell}\label{topology-d-human-on-the-loop-mission-cell}}

\begin{itemize}
\tightlist
\item
  Strength: legal and political control for high consequence actions.
\item
  Risk: operator overload and intervention lag.
\item
  ACP mitigation: consequence-tier routing and interrupt priority
  queues.
\end{itemize}

\hypertarget{d.3-swarm-mission-lifecycle-with-acp-control-gates}{%
\subsection{D.3 Swarm mission lifecycle with ACP control
gates}\label{d.3-swarm-mission-lifecycle-with-acp-control-gates}}

\hypertarget{phase-1-mission-framing}{%
\subsubsection{Phase 1: Mission framing}\label{phase-1-mission-framing}}

\begin{itemize}
\tightlist
\item
  Human command establishes CIP.
\item
  ACP pre-validates legal/policy compatibility against theater profile.
\end{itemize}

\hypertarget{phase-2-force-composition}{%
\subsubsection{Phase 2: Force
composition}\label{phase-2-force-composition}}

\begin{itemize}
\tightlist
\item
  Orchestrator allocates specialist agents and swarm units.
\item
  ACP checks each DAT/TEE against trust scope policy.
\end{itemize}

\hypertarget{phase-3-execution-and-adaptation}{%
\subsubsection{Phase 3: Execution and
adaptation}\label{phase-3-execution-and-adaptation}}

\begin{itemize}
\tightlist
\item
  Agents execute subtasks in parallel.
\item
  ACP enforces runtime policy, budget, and boundary checks per action.
\end{itemize}

\hypertarget{phase-4-degraded-operation}{%
\subsubsection{Phase 4: Degraded
operation}\label{phase-4-degraded-operation}}

\begin{itemize}
\tightlist
\item
  Under EW/comms impairment, swarms transition to preapproved degraded
  profiles.
\item
  ACP requires deterministic downgrade behavior, not emergent
  improvisation.
\end{itemize}

\hypertarget{phase-5-effect-verification}{%
\subsubsection{Phase 5: Effect
verification}\label{phase-5-effect-verification}}

\begin{itemize}
\tightlist
\item
  Mission effects must be verified through independent telemetry or
  fused confidence checks.
\item
  Unverified effects cannot be promoted as completed objectives.
\end{itemize}

\hypertarget{phase-6-termination-and-reconstitution}{%
\subsubsection{Phase 6: Termination and
reconstitution}\label{phase-6-termination-and-reconstitution}}

\begin{itemize}
\tightlist
\item
  Authority tokens expire, are renewed, or revoked.
\item
  Recovery workflows restore trusted state after fault or compromise.
\end{itemize}

\hypertarget{d.4-failure-taxonomy-for-mission-control-in-swarm-conflict}{%
\subsection{D.4 Failure taxonomy for mission control in swarm
conflict}\label{d.4-failure-taxonomy-for-mission-control-in-swarm-conflict}}

\hypertarget{f1.-intent-dilution}{%
\subsubsection{F1. Intent dilution}\label{f1.-intent-dilution}}

Subordinate agents optimize local proxy metrics and drift from commander
effect.

\begin{itemize}
\tightlist
\item
  Counter: intent anchoring checks at each task split.
\end{itemize}

\hypertarget{f2.-policy-race-conditions}{%
\subsubsection{F2. Policy race
conditions}\label{f2.-policy-race-conditions}}

Parallel agents trigger conflicting actions faster than governance
checks converge.

\begin{itemize}
\tightlist
\item
  Counter: atomic policy arbitration channels and action locks for
  shared resources.
\end{itemize}

\hypertarget{f3.-context-poisoning}{%
\subsubsection{F3. Context poisoning}\label{f3.-context-poisoning}}

Untrusted observations contaminate planner state and cascade into swarm
behavior.

\begin{itemize}
\tightlist
\item
  Counter: provenance tagging, confidence weighting, adversarial context
  filters.
\end{itemize}

\hypertarget{f4.-tool-compromise-connector-hijack}{%
\subsubsection{F4. Tool compromise / connector
hijack}\label{f4.-tool-compromise-connector-hijack}}

Trusted action channels become adversary-controlled pathways.

\begin{itemize}
\tightlist
\item
  Counter: connector attestation, runtime anomaly detection, immediate
  credential rotation.
\end{itemize}

\hypertarget{f5.-comms-partition-collapse}{%
\subsubsection{F5. Comms partition
collapse}\label{f5.-comms-partition-collapse}}

Disconnected swarm segments diverge and can no longer guarantee common
operating logic.

\begin{itemize}
\tightlist
\item
  Counter: partition-safe mission envelopes with hard stop rules.
\end{itemize}

\hypertarget{f6.-forensic-insufficiency}{%
\subsubsection{F6. Forensic
insufficiency}\label{f6.-forensic-insufficiency}}

Actions occurred, effects happened, but attribution and legal
reconstruction are impossible.

\begin{itemize}
\tightlist
\item
  Counter: mandatory EP completeness checks before mission closure.
\end{itemize}

\hypertarget{d.5-red-team-inject-library-mission-control-stress-tests}{%
\subsection{D.5 Red-team inject library (mission-control stress
tests)}\label{d.5-red-team-inject-library-mission-control-stress-tests}}

\hypertarget{inject-family-a-data-and-telemetry-deception}{%
\subsubsection{Inject family A: Data and telemetry
deception}\label{inject-family-a-data-and-telemetry-deception}}

\begin{itemize}
\tightlist
\item
  spoofed ISR targets,
\item
  delayed sensor feeds,
\item
  contradictory multi-source observations.
\end{itemize}

\hypertarget{inject-family-b-orchestration-abuse}{%
\subsubsection{Inject family B: Orchestration
abuse}\label{inject-family-b-orchestration-abuse}}

\begin{itemize}
\tightlist
\item
  fake agent identity announcements,
\item
  cross-agent impersonation,
\item
  unauthorized delegation propagation.
\end{itemize}

\hypertarget{inject-family-c-resource-exhaustion}{%
\subsubsection{Inject family C: Resource
exhaustion}\label{inject-family-c-resource-exhaustion}}

\begin{itemize}
\tightlist
\item
  decoy swarms to drain interceptors,
\item
  command queue flooding,
\item
  policy engine saturation attempts.
\end{itemize}

\hypertarget{inject-family-d-legal-boundary-pressure}{%
\subsubsection{Inject family D: Legal boundary
pressure}\label{inject-family-d-legal-boundary-pressure}}

\begin{itemize}
\tightlist
\item
  borderline target discrimination,
\item
  adversary human-shield framing,
\item
  deceptive surrender/medical signals.
\end{itemize}

\hypertarget{inject-family-e-recovery-denial}{%
\subsubsection{Inject family E: Recovery
denial}\label{inject-family-e-recovery-denial}}

\begin{itemize}
\tightlist
\item
  compromised node rejoin attempts,
\item
  stale-policy replay,
\item
  corrupted mission logs at reintegration.
\end{itemize}

\hypertarget{d.6-mission-control-readiness-model-acp-mc-levels}{%
\subsection{D.6 Mission Control readiness model (ACP-MC
levels)}\label{d.6-mission-control-readiness-model-acp-mc-levels}}

\hypertarget{level-0-ad-hoc-autonomy}{%
\subsubsection{Level 0: Ad hoc autonomy}\label{level-0-ad-hoc-autonomy}}

\begin{itemize}
\tightlist
\item
  isolated tools, no enforceable trust scopes.
\end{itemize}

\hypertarget{level-1-structured-experimentation}{%
\subsubsection{Level 1: Structured
experimentation}\label{level-1-structured-experimentation}}

\begin{itemize}
\tightlist
\item
  pilots with basic logging and manual guardrails.
\end{itemize}

\hypertarget{level-2-governed-deployment}{%
\subsubsection{Level 2: Governed
deployment}\label{level-2-governed-deployment}}

\begin{itemize}
\tightlist
\item
  trust scopes, policy checks, and role-bound execution in production.
\end{itemize}

\hypertarget{level-3-contested-ready-orchestration}{%
\subsubsection{Level 3: Contested-ready
orchestration}\label{level-3-contested-ready-orchestration}}

\begin{itemize}
\tightlist
\item
  degraded-mode doctrine + partition-safe operation + red-team
  certification.
\end{itemize}

\hypertarget{level-4-coalition-interoperable-mission-control}{%
\subsubsection{Level 4: Coalition-interoperable mission
control}\label{level-4-coalition-interoperable-mission-control}}

\begin{itemize}
\tightlist
\item
  protocol federation with cross-org policy and evidence compatibility.
\end{itemize}

\hypertarget{level-5-adaptive-strategic-control-plane}{%
\subsubsection{Level 5: Adaptive strategic control
plane}\label{level-5-adaptive-strategic-control-plane}}

\begin{itemize}
\tightlist
\item
  continuous model/tool governance, campaign-scale optimization, and
  rapid reconstitution at force level.
\end{itemize}

\hypertarget{d.7-campaign-archetypes-where-acp-mission-control-is-decisive}{%
\subsection{D.7 Campaign archetypes where ACP Mission Control is
decisive}\label{d.7-campaign-archetypes-where-acp-mission-control-is-decisive}}

\hypertarget{archetype-1-isr-to-fire-compression-campaign}{%
\subsubsection{Archetype 1: ISR-to-fire compression
campaign}\label{archetype-1-isr-to-fire-compression-campaign}}

\begin{itemize}
\tightlist
\item
  objective: reduce target cycle time while maintaining discrimination.
\item
  key ACP controls: confidence thresholds, dual-source verification,
  consequence-tier gating.
\end{itemize}

\hypertarget{archetype-2-logistics-corridor-under-drone-harassment}{%
\subsubsection{Archetype 2: Logistics corridor under drone
harassment}\label{archetype-2-logistics-corridor-under-drone-harassment}}

\begin{itemize}
\tightlist
\item
  objective: sustain throughput under persistent attrition pressure.
\item
  key ACP controls: convoy swarm overwatch orchestration, route policy
  adaptation, autonomous fallback routing.
\end{itemize}

\hypertarget{archetype-3-base-and-force-concentration-c-uas-defense}{%
\subsubsection{Archetype 3: Base and force concentration C-UAS
defense}\label{archetype-3-base-and-force-concentration-c-uas-defense}}

\begin{itemize}
\tightlist
\item
  objective: survive saturation and preserve operational continuity.
\item
  key ACP controls: priority queueing, magazine-aware engagement logic,
  layered defensive coordination.
\end{itemize}

\hypertarget{archetype-4-multi-domain-feintdeception-battlespace}{%
\subsubsection{Archetype 4: Multi-domain feint/deception
battlespace}\label{archetype-4-multi-domain-feintdeception-battlespace}}

\begin{itemize}
\tightlist
\item
  objective: avoid adversary manipulation and preserve decision quality.
\item
  key ACP controls: provenance-centric context fusion, anti-replay
  messaging, deception-resilient planning loops.
\end{itemize}

\hypertarget{d.8-acquisition-and-force-development-implications}{%
\subsection{D.8 Acquisition and force-development
implications}\label{d.8-acquisition-and-force-development-implications}}

\hypertarget{requirement-1-no-autonomy-fielding-without-acp-compatibility-profile}{%
\subsubsection{Requirement 1: No autonomy fielding without ACP
compatibility
profile}\label{requirement-1-no-autonomy-fielding-without-acp-compatibility-profile}}

Any agentic capability should publish an ACP compatibility declaration:

\begin{itemize}
\tightlist
\item
  trust scope model,
\item
  policy integration points,
\item
  evidence schema conformance,
\item
  degraded-mode behavior,
\item
  revocation mechanism.
\end{itemize}

\hypertarget{requirement-2-contractual-enforcement-clauses}{%
\subsubsection{Requirement 2: Contractual enforcement
clauses}\label{requirement-2-contractual-enforcement-clauses}}

Programs should include enforceable clauses for:

\begin{itemize}
\tightlist
\item
  security attestation of connectors and skills,
\item
  mandatory telemetry and evidence output formats,
\item
  continuous evaluation and rollback rights,
\item
  authority for emergency disablement.
\end{itemize}

\hypertarget{requirement-3-training-and-command-education}{%
\subsubsection{Requirement 3: Training and command
education}\label{requirement-3-training-and-command-education}}

Mission Control competence must become a command skill:

\begin{itemize}
\tightlist
\item
  delegation design,
\item
  autonomy risk budgeting,
\item
  forensic interpretation,
\item
  degraded-mode battle drills.
\end{itemize}

\hypertarget{d.9-working-conclusion-current}{%
\subsection{D.9 Working conclusion
(current)}\label{d.9-working-conclusion-current}}

The core strategic question is not whether swarms or agent orchestration
will matter. They already do.

The strategic question is whether DoW can establish Mission Control
discipline faster than adversaries can scale autonomous action.

ACP is therefore not merely an architecture choice. It is a warfighting
imperative tied to:

\begin{itemize}
\tightlist
\item
  tempo,
\item
  legitimacy,
\item
  survivability,
\item
  and campaign-level adaptability.
\end{itemize}

\hypertarget{d.10-open-source-intake-queue-pending-authenticated-extraction}{%
\subsection{D.10 Open source intake queue (pending authenticated
extraction)}\label{d.10-open-source-intake-queue-pending-authenticated-extraction}}

\begin{itemize}
\tightlist
\item
  X post submitted by user:
  \url{https://x.com/vraserx/status/2021654971343610058?s=52}
\item
  X post submitted by user:
  \url{https://x.com/theinformation/status/2021707053727662407?s=46}
\end{itemize}

Current status: links captured as evidence pointers; full text
extraction blocked in this environment due X authentication constraints.

\hypertarget{e.0-exhaustive-academic-paper-sweep-current-pass}{%
\subsection{E.0 Exhaustive academic paper sweep (current
pass)}\label{e.0-exhaustive-academic-paper-sweep-current-pass}}

Goal: create a broad, high-signal academic corpus spanning MARL, swarms,
agent orchestration, safety/governance, and adversarial resilience for
ACP doctrine development.

Method (this pass):

\begin{itemize}
\tightlist
\item
  Queried OpenAlex across focused topic clusters.
\item
  Applied relevance filters and citation thresholds.
\item
  Added canonical papers manually where needed.
\item
  Preserved links for direct retrieval.
\end{itemize}

Notes:

\begin{itemize}
\tightlist
\item
  This is intentionally broad to maximize coverage for follow-on
  pruning.
\item
  Some cross-domain papers are retained because methods transfer
  directly to military autonomy contexts (coordination, robustness,
  safety constraints, distributed control).
\end{itemize}

\hypertarget{marl-core}{%
\subsubsection{MARL core}\label{marl-core}}

\begin{enumerate}
\def\labelenumi{\arabic{enumi}.}
\tightlist
\item
  Multi-Agent Reinforcement Learning: Independent vs.~Cooperative Agents
  (1993), cited by 1789.
  \url{https://doi.org/10.1016/b978-1-55860-307-3.50049-6}
\item
  Counterfactual Multi-Agent Policy Gradients (2018), cited by 1543.
  \url{https://doi.org/10.1609/aaai.v32i1.11794}
\item
  Multi-Agent Actor-Critic for Mixed Cooperative-Competitive
  Environments (2017), cited by 1014.
  \url{http://arxiv.org/abs/1706.02275}
\item
  Cooperative Multi-agent Control Using Deep Reinforcement Learning
  (2017), cited by 846.
  \url{https://doi.org/10.1007/978-3-319-71682-4_5}
\item
  Multi-agent deep reinforcement learning: a survey (2021), cited by
  726. \url{https://doi.org/10.1007/s10462-021-09996-w}
\item
  QMIX: Monotonic Value Function Factorisation for Deep Multi-Agent
  Reinforcement Learning (2018), cited by 475.
  \url{http://arxiv.org/abs/1803.11485}
\item
  Monotonic Value Function Factorisation for Deep Multi-Agent
  Reinforcement Learning (2020), cited by 429.
  \url{http://arxiv.org/abs/2003.08839}
\item
  A review of cooperative multi-agent deep reinforcement learning
  (2022), cited by 404. \url{https://doi.org/10.1007/s10489-022-04105-y}
\item
  QMIX: Monotonic Value Function Factorisation for Deep Multi-Agent
  Reinforcement Learning (2018), cited by 352.
  \url{http://arxiv.org/abs/1803.11485}
\item
  Multi-Agent Reinforcement Learning: A Review of Challenges and
  Applications (2021), cited by 315.
  \url{https://doi.org/10.3390/app11114948}
\item
  A novel optimal bipartite consensus control scheme for unknown
  multi-agent systems via model-free reinforcement learning (2019),
  cited by 298. \url{https://doi.org/10.1016/j.amc.2019.124821}
\item
  Actor-Attention-Critic for Multi-Agent Reinforcement Learning (2018),
  cited by 289. \url{http://arxiv.org/abs/1810.02912}
\item
  Robust Multi-Agent Reinforcement Learning via Minimax Deep
  Deterministic Policy Gradient (2019), cited by 279.
  \url{https://doi.org/10.1609/aaai.v33i01.33014213}
\item
  A survey on multi-agent deep reinforcement learning: from the
  perspective of challenges and applications (2020), cited by 237.
  \url{https://doi.org/10.1007/s10462-020-09938-y}
\item
  QTRAN: Learning to Factorize with Transformation for Cooperative
  Multi-Agent Reinforcement Learning (2019), cited by 222.
  \url{https://doi.org/10.48550/arxiv.1905.05408}
\item
  QTRAN: Learning to Factorize with Transformation for Cooperative
  Multi-Agent Reinforcement Learning (2019), cited by 198.
  \url{http://arxiv.org/abs/1905.05408}
\item
  Multi-Agent Deep Reinforcement Learning for Multi-Robot Applications:
  A Survey (2023), cited by 140. \url{https://doi.org/10.3390/s23073625}
\item
  Dealing with Non-Stationarity in Multi-Agent Deep Reinforcement
  Learning (2019), cited by 124. \url{http://arxiv.org/abs/1906.04737}
\item
  Evolutionary game theory and multi-agent reinforcement learning
  (2005), cited by 123. \url{https://doi.org/10.1017/s026988890500041x}
\item
  Multi-Agent Actor-Critic with Hierarchical Graph Attention Network
  (2020), cited by 118. \url{https://doi.org/10.1609/aaai.v34i05.6214}
\item
  Value-Decomposition Networks For Cooperative Multi-Agent Learning
  (2017). \url{https://arxiv.org/abs/1706.05296}
\item
  Learning to Communicate with Deep Multi-Agent Reinforcement Learning
  (2016). \url{https://arxiv.org/abs/1605.06676}
\end{enumerate}

\hypertarget{swarm-robotics-hsi}{%
\subsubsection{Swarm robotics + HSI}\label{swarm-robotics-hsi}}

\begin{enumerate}
\def\labelenumi{\arabic{enumi}.}
\tightlist
\item
  Swarm robotics: a review from the swarm engineering perspective
  (2013), cited by 1638. \url{https://doi.org/10.1007/s11721-012-0075-2}
\item
  A Survey on Aerial Swarm Robotics (2018), cited by 604.
  \url{https://doi.org/10.1109/tro.2018.2857475}
\item
  Swarm of micro flying robots in the wild (2022), cited by 490.
  \url{https://doi.org/10.1126/scirobotics.abm5954}
\item
  A review of swarm robotics tasks (2015), cited by 428.
  \url{https://doi.org/10.1016/j.neucom.2015.05.116}
\item
  Swarm Robotic Behaviors and Current Applications (2020), cited by 413.
  \url{https://doi.org/10.3389/frobt.2020.00036}
\item
  Human Interaction With Robot Swarms: A Survey (2015), cited by 387.
  \url{https://doi.org/10.1109/thms.2015.2480801}
\item
  Swarm Robotics: Past, Present, and Future {[}Point of View{]} (2021),
  cited by 327. \url{https://doi.org/10.1109/jproc.2021.3072740}
\item
  Optimized Stochastic Policies for Task Allocation in Swarms of Robots
  (2009), cited by 244. \url{https://doi.org/10.1109/tro.2009.2024997}
\item
  Search and tracking algorithms for swarms of robots: A survey (2015),
  cited by 180. \url{https://doi.org/10.1016/j.robot.2015.08.010}
\item
  Predictive control of aerial swarms in cluttered environments (2021),
  cited by 151. \url{https://doi.org/10.1038/s42256-021-00341-y}
\item
  A Survey on Robotic Technologies for Forest Firefighting: Applying
  Drone Swarms to Improve Firefighters' Efficiency and Safety (2021),
  cited by 134. \url{https://doi.org/10.3390/app11010363}
\item
  Speaking Swarmish: Human-Robot Interface Design for Large Swarms of
  Autonomous Mobile Robots. (2006), cited by 113.
  \url{http://citeseerx.ist.psu.edu/viewdoc/summary?doi=10.1.1.130.9920}
\item
  Human swarm interaction for radiation source search and localization
  (2008), cited by 112. \url{https://doi.org/10.1109/sis.2008.4668287}
\item
  Human Swarm Interaction: An Experimental Study of Two Types of
  Interaction with Foraging Swarms (2013), cited by 104.
  \url{https://doi.org/10.5898/jhri.2.2.kolling}
\item
  Mean-field models in swarm robotics: a survey (2019), cited by 94.
  \url{https://doi.org/10.1088/1748-3190/ab49a4}
\item
  Aerial Shepherds: Coordination among UAVs and Swarms of Robots (2008),
  cited by 92. \url{https://doi.org/10.1007/978-4-431-35873-2_24}
\item
  Swarm robots in mechanized agricultural operations: A review about
  challenges for research (2022), cited by 91.
  \url{https://doi.org/10.1016/j.compag.2021.106608}
\item
  From animal collective behaviors to swarm robotic cooperation (2023),
  cited by 86. \url{https://doi.org/10.1093/nsr/nwad040}
\item
  Benchmark of swarm robotics distributed techniques in a search task
  (2013), cited by 86. \url{https://doi.org/10.1016/j.robot.2013.10.004}
\item
  Human-swarm interaction using spatial gestures (2014), cited by 85.
  \url{https://doi.org/10.1109/iros.2014.6943101}
\end{enumerate}

\hypertarget{llm-agents-orchestration}{%
\subsubsection{LLM agents +
orchestration}\label{llm-agents-orchestration}}

\begin{enumerate}
\def\labelenumi{\arabic{enumi}.}
\tightlist
\item
  A survey on large language model based autonomous agents (2024), cited
  by 797. \url{https://doi.org/10.1007/s11704-024-40231-1}
\item
  A survey on LLM-based multi-agent systems: workflow, infrastructure,
  and challenges (2024), cited by 153.
  \url{https://doi.org/10.1007/s44336-024-00009-2}
\item
  AutoGen: Enabling Next-Gen LLM Applications via Multi-Agent
  Conversation (2023), cited by 138.
  \url{http://arxiv.org/abs/2308.08155}
\item
  Orchestrated Scheduling and Multi-Agent Deep Reinforcement Learning
  for Cloud-Assisted Multi-UAV Charging Systems (2021), cited by 93.
  \url{https://doi.org/10.1109/tvt.2021.3062418}
\item
  Towards autonomous system: flexible modular production system enhanced
  with large language model agents (2023), cited by 64.
  \url{https://doi.org/10.1109/etfa54631.2023.10275362}
\item
  A multi-agent and cloud-edge orchestration framework of digital twin
  for distributed production control (2023), cited by 62.
  \url{https://doi.org/10.1016/j.rcim.2023.102543}
\item
  ChatEval: Towards Better LLM-based Evaluators through Multi-Agent
  Debate (2023), cited by 51. \url{http://arxiv.org/abs/2308.07201}
\item
  Large Language Model based Multi-Agents: A Survey of Progress and
  Challenges (2024), cited by 50. \url{http://arxiv.org/abs/2402.01680}
\item
  Stance Detection with Collaborative Role-Infused LLM-Based Agents
  (2024), cited by 44. \url{https://doi.org/10.1609/icwsm.v18i1.31360}
\item
  Exploring Large Language Model based Intelligent Agents: Definitions,
  Methods, and Prospects (2024), cited by 35.
  \url{http://arxiv.org/abs/2401.03428}
\item
  Personal LLM Agents: Insights and Survey about the Capability,
  Efficiency and Security (2024), cited by 30.
  \url{http://arxiv.org/abs/2401.05459}
\item
  Understanding the planning of LLM agents: A survey (2024), cited by
  26. \url{http://arxiv.org/abs/2402.02716}
\item
  Encouraging Divergent Thinking in Large Language Models through
  Multi-Agent Debate (2023), cited by 26.
  \url{http://arxiv.org/abs/2305.19118}
\item
  From LLMs to LLM-based Agents for Software Engineering: A Survey of
  Current, Challenges and Future (2024), cited by 15.
  \url{http://arxiv.org/abs/2408.02479}
\item
  PRefLexOR: preference-based recursive language modeling for
  exploratory optimization of reasoning and agentic thinking (2025),
  cited by 13. \url{https://doi.org/10.1038/s44387-025-00003-z}
\item
  LLM-based Multi-Agent Reinforcement Learning: Current and Future
  Directions (2024), cited by 10. \url{http://arxiv.org/abs/2405.11106}
\item
  Thematic-LM: A LLM-based Multi-agent System for Large-scale Thematic
  Analysis (2025), cited by 8.
  \url{https://doi.org/10.1145/3696410.3714595}
\item
  LLM-Based Multi-Agent Systems for Software Engineering: Literature
  Review, Vision and the Road Ahead (2024), cited by 8.
  \url{http://arxiv.org/abs/2404.04834}
\item
  Towards large language model-based personal agents in the enterprise:
  Current trends and open problems (2023), cited by 8.
  \url{https://doi.org/10.18653/v1/2023.findings-emnlp.461}
\item
  TradingGPT: Multi-Agent System with Layered Memory and Distinct
  Characters for Enhanced Financial Trading Performance (2023), cited by
  8. \url{http://arxiv.org/abs/2309.03736}
\item
  SWE-agent: Agent-Computer Interfaces Enable Automated Software
  Engineering (2024). \url{https://arxiv.org/abs/2405.15793}
\item
  CAMEL: Communicative Agents for Mind Exploration of Large Language
  Model Society (2023). \url{https://arxiv.org/abs/2303.17760}
\item
  AgentVerse: Facilitating Multi-Agent Collaboration and Exploring
  Emergent Behaviors (2023). \url{https://arxiv.org/abs/2308.10848}
\item
  MetaGPT: Meta Programming for Multi-Agent Collaborative Framework
  (2023). \url{https://arxiv.org/abs/2308.00352}
\end{enumerate}

\hypertarget{safetygovernancelaw}{%
\subsubsection{Safety/governance/LAW}\label{safetygovernancelaw}}

\begin{enumerate}
\def\labelenumi{\arabic{enumi}.}
\tightlist
\item
  Reinforcement Learning: A Survey (1996), cited by 8635.
  \url{https://doi.org/10.1613/jair.301}
\item
  A Multiagent Approach to Autonomous Intersection Management (2008),
  cited by 1283. \url{https://doi.org/10.1613/jair.2502}
\item
  A comprehensive survey on safe reinforcement learning (2015), cited by
  1191.
  \url{https://jmlr.csail.mit.edu/papers/volume16/garcia15a/garcia15a.pdf}
\item
  Continuous control for robot based on deep reinforcement learning
  (2019), cited by 854. \url{https://doi.org/10.32657/10356/90191}
\item
  Safe Learning in Robotics: From Learning-Based Control to Safe
  Reinforcement Learning (2022), cited by 594.
  \url{https://doi.org/10.1146/annurev-control-042920-020211}
\item
  A Tour of Reinforcement Learning: The View from Continuous Control
  (2018), cited by 573.
  \url{https://doi.org/10.1146/annurev-control-053018-023825}
\item
  Trustworthy artificial intelligence (2020), cited by 507.
  \url{https://doi.org/10.1007/s12525-020-00441-4}
\item
  Meaningful Human Control over Autonomous Systems: A Philosophical
  Account (2018), cited by 405.
  \url{https://doi.org/10.3389/frobt.2018.00015}
\item
  On banning autonomous weapon systems: human rights, automation, and
  the dehumanization of lethal decision-making (2012), cited by 339.
  \url{https://doi.org/10.1017/s1816383112000768}
\item
  Reinforcement learning in robotic applications: a comprehensive survey
  (2021), cited by 318. \url{https://doi.org/10.1007/s10462-021-09997-9}
\item
  A Review of Trustworthy and Explainable Artificial Intelligence (XAI)
  (2023), cited by 225.
  \url{https://doi.org/10.1109/access.2023.3294569}
\item
  Improving Robot Controller Transparency Through Autonomous Policy
  Explanation (2017), cited by 219.
  \url{https://doi.org/10.1145/2909824.3020233}
\item
  Responsibility Research for Trustworthy Autonomous Systems (Blue Sky
  Ideas Track) (2021), cited by 218.
  \url{https://eprints.soton.ac.uk/447511/1/Responsibility_Research_for_Trustworthy_Autonomous_Systems.pdf}
\item
  Governance Strategies for a Sustainable Digital World (2018), cited by
  215. \url{https://doi.org/10.3390/su10020440}
\item
  Autonomous weapons systems, killer robots and human dignity (2018),
  cited by 176. \url{https://doi.org/10.1007/s10676-018-9494-0}
\item
  Overcoming Barriers to Cross-cultural Cooperation in AI Ethics and
  Governance (2020), cited by 160.
  \url{https://doi.org/10.1007/s13347-020-00402-x}
\item
  The Law of Armed Conflict: International Humanitarian Law in War
  (2021), cited by 124.
  \url{http://bvbr.bib-bvb.de:8991/F?func=service\&doc_library=BVB01\&local_base=BVB01\&doc_number=020516470\&sequence=000001\&line_number=0001\&func_code=DB_RECORDS\&service_type=MEDIA}
\item
  Autonomous Weapons and International Humanitarian Law: Advantages,
  Open Technical Questions and Legal Issues to be Clarified (2014),
  cited by 113. \url{https://archive-ouverte.unige.ch/unige:37976}
\item
  Autonomous Weapon Systems and International Humanitarian Law: A Reply
  to the Critics (2012), cited by 101.
  \url{https://doi.org/10.2139/ssrn.2184826}
\item
  Open Problems and Fundamental Limitations of Reinforcement Learning
  from Human Feedback (2023), cited by 88.
  \url{http://arxiv.org/abs/2307.15217}
\end{enumerate}

\hypertarget{adversarial-resilience}{%
\subsubsection{Adversarial resilience}\label{adversarial-resilience}}

\begin{enumerate}
\def\labelenumi{\arabic{enumi}.}
\tightlist
\item
  Secure impulsive synchronization control of multi-agent systems under
  deception attacks (2018), cited by 324.
  \url{https://doi.org/10.1016/j.ins.2018.04.020}
\item
  Secure impulsive synchronization in Lipschitz-type multi-agent systems
  subject to deception attacks (2020), cited by 165.
  \url{https://doi.org/10.1109/jas.2020.1003297}
\item
  A study on cyber-security of autonomous and unmanned vehicles (2015),
  cited by 136. \url{https://doi.org/10.1177/1548512915575803}
\item
  An analysis of the robustness of UAV agriculture field coverage using
  multi-agent reinforcement learning (2023), cited by 124.
  \url{https://doi.org/10.1007/s41870-023-01264-0}
\item
  Cooperative adaptive fault-tolerant control for multi-agent systems
  with deception attacks (2020), cited by 111.
  \url{https://doi.org/10.1016/j.jfranklin.2019.12.032}
\item
  Detection of Cyber-attacks to indoor real time localization systems
  for autonomous robots (2017), cited by 87.
  \url{https://doi.org/10.1016/j.robot.2017.10.006}
\item
  Fault-Tolerant secure consensus tracking of delayed nonlinear
  multi-agent systems with deception attacks and uncertain parameters
  via impulsive control (2019), cited by 85.
  \url{https://doi.org/10.1016/j.cnsns.2019.105043}
\item
  Robust Regional Coordination of Inverter-Based Volt/Var Control via
  Multi-Agent Deep Reinforcement Learning (2021), cited by 70.
  \url{https://doi.org/10.1109/tsg.2021.3104139}
\item
  Enhancing Autonomous System Security and Resilience With Generative
  AI: A Comprehensive Survey (2024), cited by 60.
  \url{https://doi.org/10.1109/access.2024.3439363}
\item
  Failure-Scenario Maker for Rule-Based Agent using Multi-agent
  Adversarial Reinforcement Learning and its Application to Autonomous
  Driving (2019), cited by 57.
  \url{https://doi.org/10.24963/ijcai.2019/832}
\item
  Optimal Bi-Level Bidding and Dispatching Strategy Between Active
  Distribution Network and Virtual Alliances Using Distributed Robust
  Multi-Agent Deep Reinforcement Learning (2022), cited by 56.
  \url{https://doi.org/10.1109/tsg.2022.3164080}
\item
  Securing unmanned autonomous systems from cyber threats (2016), cited
  by 50. \url{https://doi.org/10.1177/1548512916628335}
\item
  Non-fragile consensus control for nonlinear multi-agent systems with
  uniform quantizations and deception attacks via output feedback
  approach (2019), cited by 49.
  \url{https://doi.org/10.1007/s11071-019-04787-z}
\item
  Towards coordinated and robust real-time control: a decentralized
  approach for combined sewer overflow and urban flooding reduction
  based on multi-agent reinforcement learning (2022), cited by 48.
  \url{https://doi.org/10.1016/j.watres.2022.119498}
\item
  Multi-Agent Adversarial Inverse Reinforcement Learning (2019), cited
  by 45. \url{http://arxiv.org/abs/1907.13220}
\item
  An asymmetric Lyapunov-Krasovskii functional approach for
  event-triggered consensus of multi-agent systems with deception
  attacks (2022), cited by 44.
  \url{https://doi.org/10.1016/j.amc.2022.127584}
\item
  Resilient finite-time distributed event-triggered consensus of
  multi-agent systems with multiple cyber-attacks (2022), cited by 44.
  \url{https://doi.org/10.1016/j.cnsns.2022.106876}
\item
  Observer-based non-fragile H∞-consensus control for multi-agent
  systems under deception attacks (2021), cited by 44.
  \url{https://doi.org/10.1080/00207721.2021.1884917}
\item
  Behaviour-Based Anomaly Detection of Cyber-Physical Attacks on a
  Robotic Vehicle (2016), cited by 44.
  \url{https://doi.org/10.1109/iucc-css.2016.017}
\item
  Resilience enhancement of multi-agent reinforcement learning-based
  demand response against adversarial attacks (2022), cited by 43.
  \url{https://doi.org/10.1016/j.apenergy.2022.119688}
\item
  On the Robustness of Cooperative Multi-Agent Reinforcement Learning
  (2020). \url{https://arxiv.org/abs/2006.07538}
\item
  Robust Multi-Agent Reinforcement Learning via Minimax Deep
  Deterministic Policy Gradient (2019).
  \url{https://arxiv.org/abs/1810.11960}
\end{enumerate}

\hypertarget{e.1-optional-next-pass-v5.1-if-desired}{%
\subsection{E.1 Optional next pass (v5.1) if
desired}\label{e.1-optional-next-pass-v5.1-if-desired}}

\begin{itemize}
\tightlist
\item
  Add per-paper relevance tags: \texttt{core}, \texttt{supporting},
  \texttt{adjacent}, \texttt{drop}.
\item
  Compute citation graph clusters to identify canonical hubs.
\item
  Produce a ``DoW mission-control reading list'' shortlist (25 papers)
  for direct memo integration.
\item
  Add a ``warfighting translation'' column for each paper (what control
  primitive it informs).
\end{itemize}

\hypertarget{additional-references-for-acp-warfighting-imperative}{%
\subsection{Additional references for ACP warfighting
imperative}\label{additional-references-for-acp-warfighting-imperative}}

\begin{enumerate}
\def\labelenumi{\arabic{enumi}.}
\setcounter{enumi}{34}
\tightlist
\item
  \textbf{GAO} --- Defense Command and Control: Further Progress Hinges
  on Establishing a Comprehensive Framework (Apr 2025)
  \url{https://files.gao.gov/reports/GAO-25-106454/index.html}
\item
  \textbf{CSIS} --- Ukraine's Future Vision and Current Capabilities for
  Waging AI-Enabled Autonomous Warfare (Mar 2025)
  \url{https://www.csis.org/analysis/ukraines-future-vision-and-current-capabilities-waging-ai-enabled-autonomous-warfare}
\item
  \textbf{CSIS} --- How Russia Is Reshaping Command and Control for
  AI-Enabled Warfare (Feb 2026)
  \url{https://www.csis.org/analysis/how-russia-reshaping-command-and-control-ai-enabled-warfare}
\item
  \textbf{RAND} --- How Could Artificial Intelligence Shape the Future
  of War? (2026)
  \url{https://www.rand.org/pubs/research_reports/RRA4316-1.html}
\item
  \textbf{RAND} --- Strategic competition in the age of AI: Emerging
  risks and opportunities from military use of AI (2024)
  \url{https://www.rand.org/pubs/research_reports/RRA3295-1.html}
\item
  \textbf{NATO} --- Summary of NATO's revised AI strategy (Jul 2024)
  \url{https://www.nato.int/en/about-us/official-texts-and-resources/official-texts/2024/07/10/summary-of-natos-revised-artificial-intelligence-ai-strategy}
\item
  \textbf{UK MOD (GOV.UK)} --- JSP 936: Dependable Artificial
  Intelligence in defence (Part 1 directive) (Nov 2024)
  \url{https://www.gov.uk/government/publications/jsp-936-dependable-artificial-intelligence-ai-in-defence-part-1-directive}
\item
  \textbf{UK MOD (GOV.UK)} --- Laying the Groundwork: Responsible AI
  Senior Officers' Report 2025 (Oct 2025)
  \url{https://www.gov.uk/government/publications/laying-the-groundwork-responsible-ai-senior-officers-report-2025}
\item
  \textbf{U.S. Department of State} --- Political Declaration on
  Responsible Military Use of Artificial Intelligence and Autonomy
  (updated 2024 context)
  \url{https://www.state.gov/bureau-of-arms-control-deterrence-and-stability/political-declaration-on-responsible-military-use-of-artificial-intelligence-and-autonomy}
\item
  \textbf{UNODA} --- Artificial intelligence in the military domain
  (including UNGA resolution 79/239 context)
  \url{https://disarmament.unoda.org/en/our-work/emerging-challenges/artificial-intelligence-military-domain}
\item
  \textbf{SIPRI} --- Retaining Human Responsibility in the Development
  and Use of Autonomous Weapon Systems (Oct 2022)
  \url{https://www.sipri.org/publications/2022/policy-reports/retaining-human-responsibility-development-and-use-autonomous-weapon-systems-accountability}
\item
  \textbf{ICRC Law \& Policy Blog} --- Three lessons on AWS regulation
  to ensure accountability for IHL violations (Mar 2023)
  \url{https://blogs.icrc.org/law-and-policy/2023/03/02/three-lessons-autonomous-weapons-systems-ihl/}
\item
  \textbf{DIU} --- The Replicator Initiative (updated 2025)
  \url{https://www.diu.mil/replicator}
\item
  \textbf{Arms Control Association} --- AI and Nuclear Command and
  Control: It's Even More Complicated Than You Think (Sep 2025)
  \url{https://www.armscontrol.org/act/2025-09/features/artificial-intelligence-and-nuclear-command-and-control-its-even-more}
\item
  \textbf{Microsoft Copilot Blog} --- Multi-agent orchestration
  announcements at Build 2025
  \url{https://www.microsoft.com/en-us/microsoft-copilot/blog/copilot-studio/multi-agent-orchestration-maker-controls-and-more-microsoft-copilot-studio-announcements-at-microsoft-build-2025/}
\item
  \textbf{Google Cloud Blog} --- Build and manage multi-system agents
  with Vertex AI (Apr 2025)
  \url{https://cloud.google.com/blog/products/ai-machine-learning/build-and-manage-multi-system-agents-with-vertex-ai}
\item
  \textbf{Google Cloud Blog} --- Agent2Agent protocol (A2A) is getting
  an upgrade (Jul 2025)
  \url{https://cloud.google.com/blog/products/ai-machine-learning/agent2agent-protocol-is-getting-an-upgrade}
\item
  \textbf{Google Developers Blog} --- Google Cloud donates A2A to Linux
  Foundation (Jun 2025)
  \url{https://developers.googleblog.com/en/google-cloud-donates-a2a-to-linux-foundation/}
\item
  \textbf{U.S. Army} --- Swarm Technology in Sustainment Operations (Jan
  2025)
  \url{https://www.army.mil/article/282467/swarm_technology_in_sustainment_operations}
\item
  \textbf{CNAS} --- Countering the Swarm: Protecting the Joint Force in
  the Drone Age (Sep 2025)
  \url{https://s3.us-east-1.amazonaws.com/files.cnas.org/documents/Report_CUAS_Defense_Sep-2025_final.pdf}
\item
  \textbf{CNA} --- PRC Concepts for UAV Swarms in Future Warfare (Oct
  2025)
  \url{https://www.cna.org/analyses/2025/10/prc-concepts-for-uav-swarms-in-future-warfare}
\item
  \textbf{U.S. Secretary of Defense Memorandum} --- Replicator 2
  Direction and Execution (Sep 27, 2024)
  \url{https://s3.us-gov-west-1.amazonaws.com/assets.diu.mil/3nanhbfkr0pc/1dkJGhMeAgPldz1nnIwabK/abf85531a4281cddab6b0d8c953440e2/REPLICATOR-2-MEMO-SD-SIGNED__1_.pdf}
\item
  \textbf{X source pointer (The Information)} --- post link queued for
  authenticated extraction
  \url{https://x.com/theinformation/status/2021707053727662407?s=46}
\end{enumerate}

\hypertarget{build-target-for-review}{%
\subsection{Build target for review}\label{build-target-for-review}}

\begin{itemize}
\tightlist
\item
  Generated \texttt{.tex} source path (internal):
  \texttt{Whitepaper/writing/notes/tex/2026-02-12-annex-a-supporting-notes-v5.tex}
\item
  Generated PDF artifact path (internal):
  \texttt{out/memos/2026-02-12-annex-a-supporting-notes-v5.pdf}
\item
  Publication status: \textbf{draft / non-published}.
\end{itemize}

\end{document}
